
\documentclass{article}
\usepackage{geometry}
\geometry{margin=1.5cm, vmargin={0pt,1cm}}
\setlength{\topmargin}{-1cm}
\setlength{\headheight}{12.5pt}
\setlength{\paperheight}{29.7cm}
\setlength{\textheight}{25.3cm}

% useful packages.
\usepackage[UTF8]{ctex}
\usepackage{fancyhdr}
\usepackage{layout}
\usepackage{luacode}
\usepackage{listings}
\usepackage{xcolor}

% Configure listings for Python code highlighting
\lstset{
  language=Python,
  basicstyle=\ttfamily\small,
  keywordstyle=\bfseries\color{blue},
  commentstyle=\color{gray},
  stringstyle=\color{red},
  breaklines=true,
  showstringspaces=false,
  tabsize=4,
  frame=single,
  framesep=5pt,
  rulecolor=\color{gray},
  backgroundcolor=\color{white},
  captionpos=b,
  numberstyle=\tiny\color{gray},
  numbers=left,
  stepnumber=5,
}

\lstnewenvironment{plaintext}[1][]{
  \lstset{
    language=none,
    basicstyle=\ttfamily,
    keywordstyle=,
    commentstyle=,
    stringstyle=,
    numbers=none,
    frame=none,
    backgroundcolor=\color{white},
    #1                % 允许用户自行覆盖默认设置
  }
}{}

% import common macros
% % % % % % % % % % % % % % % % % % % % % % % % % % % % % % % %
% Common Macros for LaTeX
% % % % % % % % % % % % % % % % % % % % % % % % % % % % % % % %

% Useful packages
\usepackage{amsfonts}
\usepackage{amsmath}
\usepackage{amssymb}
\usepackage{amsthm}
\usepackage{enumerate}
\usepackage{graphicx}
\usepackage{multicol}
\usepackage[ruled,linesnumbered]{algorithm2e}

% Math operators and common symbols
\newcommand{\dif}{\mathrm{d}}
\newcommand{\innerProd}[1]{\left\langle #1 \right\rangle}
\newcommand{\difFrac}[2]{\frac{\dif #1}{\dif #2}}
\newcommand{\pdfFrac}[2]{\frac{\partial #1}{\partial #2}}
\newcommand{\T}{\mathrm{T}}
\newcommand{\norm}[1]{\left\| #1 \right\|}
\newcommand{\abs}[1]{\left| #1 \right|}

% Vector and Matrix shortcuts
\newcommand{\vecTwo}[2]{
  \begin{bmatrix}
    #1 \\
    #2
\end{bmatrix}}
\newcommand{\vecThree}[3]{
  \begin{bmatrix}
    #1 \\
    #2 \\
    #3
\end{bmatrix}}
\newcommand{\vecTwoH}[2]{
  \begin{bmatrix}
    #1 & #2
  \end{bmatrix}
}
\newcommand{\vecThreeH}[3]{
  \begin{bmatrix}
    #1 & #2 & #3
  \end{bmatrix}
}

% Common fields
\newcommand{\R}{\mathbb{R}}
\newcommand{\C}{\mathbb{C}}
\newcommand{\Z}{\mathbb{Z}}
\newcommand{\N}{\mathbb{N}}

% Solutions and proofs
\newenvironment{solution}{
\begin{proof}[解]}{
\end{proof}}

% Utility to get environment variables (requires lualatex)
\newcommand{\getEnv}[2]{\directlua{tex.print(os.getenv("#1") or "#2")}}


% Name and uID
\newcommand{\name}{\getEnv{NAME}{NAME}}
\newcommand{\uid}{\getEnv{UID}{UID}}
\newcommand{\course}{数值代数}
\newcommand{\hwNum}{3}

\begin{document}

\pagestyle{fancy}
\fancyhead{}
\lhead{\name (\uid)}
\chead{\course \#\hwNum}
\rhead{\today}

\section{题目 1：平方根方法}

用平方根方法求解线性方程组 $Ax = b$。

系数矩阵 $A$ 与右端向量 $b$：
\[
  A =
  \begin{pmatrix}
    4 & -2 & 4 & 2 \\
    -2 & 10 & -2 & -7 \\
    4 & -2 & 8 & 4 \\
    2 & -7 & 4 & 7
  \end{pmatrix}, \quad
  b =
  \begin{pmatrix}
    8 \\ 2 \\ 16 \\ 6
  \end{pmatrix}
\]

\begin{solution}
  计算得到 $A$ 的 Cholesky 分解 $A = LL^T$，其中
  \[
    L =
    \begin{pmatrix}
      2 & 0 & 0 & 0\\
      -1 & 3 & 0 & 0\\
      2 & 0 & 2 & 0\\
      1 & -2 & 1 & 1
    \end{pmatrix}
  \]
  从而 $Ly = b$ 的解为
  \[
    y =
    \begin{pmatrix}
      4\\2\\4\\2
    \end{pmatrix}
  \]
  最后 $L^T x = y$ 的解为

  \[
    x =
    \begin{pmatrix}
      1\\2\\1\\2
    \end{pmatrix}
  \]
\end{solution}

\section{题目 2：改进的平方根方法}

用改进的平方根方法求解线性方程组 $Ax = b$。

系数矩阵 $A$ 与右端向量 $b$：
\[
  A =
  \begin{pmatrix}
    10 & 20 & 30 \\
    20 & 45 & 80 \\
    30 & 80 & 171
  \end{pmatrix}, \quad
  b =
  \begin{pmatrix}
    10 \\ 5 \\ -31
  \end{pmatrix}
\]

\begin{solution}
  计算得到 $A$ 的改进的 Cholesky 分解 $A = LDL^T$，其中
  \[
    L =
    \begin{pmatrix}
      1 & 0 & 0\\
      2 & 1 & 0\\
      3 & 4 & 1
    \end{pmatrix}, \quad
    D =
    \begin{pmatrix}
      10 & 0 & 0\\
      0 & 5 & 0\\
      0 & 0 & 1
    \end{pmatrix}
  \]

  从而 $Ly = b$ 的解为
  \[
    y =
    \begin{pmatrix}
      10\\-15\\-1
    \end{pmatrix}
  \]

  最后 $DL^T x = y$ 的解为
  \[
    x =
    \begin{pmatrix}
      2\\1\\-1
    \end{pmatrix}
  \]

\end{solution}

\section{题目 3：追赶法求解三对角方程组}

用追赶法求解 $Ax = b$。

系数矩阵 $A$ 与右端向量 $b$：
\[
  A =
  \begin{pmatrix}
    -4 & 4 & 0 \\
    2 & -6 & 4 \\
    0 & 2 & 6
  \end{pmatrix}, \quad
  b =
  \begin{pmatrix}
    -8 \\ 8 \\ -2
  \end{pmatrix}
\]

\begin{solution}
  \[
    \begin{cases}
      \gamma_i = c_i \\
      \alpha_i = a_i - \gamma_i \beta_{i-1} \\
      \beta_i = b_i / \alpha_i
    \end{cases}
  \]
  代入计算容易得到 $A$ 的 LU 分解，其中
  \[
    L =
    \begin{pmatrix}
      -4 & & \\
      2 & -4 & \\
      & 2 & 8
    \end{pmatrix}, \quad
    U =
    \begin{pmatrix}
      1 & -1 & \\
      & 1 & -1 \\
      & & 1
    \end{pmatrix}
  \]
  从而 $Ly = b$ 的解为
  \[
    y =
    \begin{pmatrix}
      2 \\ -1 \\ \frac{1}{2}
    \end{pmatrix}
  \]
  最后 $Ux = y$ 的解为
  \[
    x =
    \begin{pmatrix}
      \frac{3}{2} \\ -\frac{1}{2} \\ \frac{1}{2}
    \end{pmatrix}
  \]
\end{solution}

\section{题目 4：Gauss-Jordan 法求矩阵逆}

用 Gauss-Jordan 法求下列矩阵 $A$ 的逆矩阵 $A^{-1}$：

\[
  A =
  \begin{pmatrix}
    2 & 1 & -3 & -1 \\
    3 & 1 & 0 & 7 \\
    -1 & 2 & 4 & -2 \\
    1 & 0 & -1 & 5
  \end{pmatrix}
\]

\begin{solution}
  用无主元 Gauss-Jordan 消元法求解。构造增广矩阵 $[A | I]$ 并通过行变换化为 $[I | A^{-1}]$。

  \textbf{初始增广矩阵 $[A | I]$：}
  \[
    \left[
      \begin{array}{cccc|cccc}
        2 & 1 & -3 & -1 & 1 & 0 & 0 & 0 \\
        3 & 1 & 0 & 7 & 0 & 1 & 0 & 0 \\
        -1 & 2 & 4 & -2 & 0 & 0 & 1 & 0 \\
        1 & 0 & -1 & 5 & 0 & 0 & 0 & 1
    \end{array}\right]
  \]

  \textbf{第 1 步：处理第 1 列}

  主元为 $2$，将第 1 行除以 2，然后消去其他行：
  \[
    \left[
      \begin{array}{cccc|cccc}
        1 & 1/2 & -3/2 & -1/2 & 1/2 & 0 & 0 & 0 \\
        0 & -1/2 & 9/2 & 17/2 & -3/2 & 1 & 0 & 0 \\
        0 & 5/2 & 5/2 & -5/2 & 1/2 & 0 & 1 & 0 \\
        0 & -1/2 & 1/2 & 11/2 & -1/2 & 0 & 0 & 1
    \end{array}\right]
  \]

  \textbf{第 2 步：处理第 2 列}

  主元为 $-1/2$，将第 2 行除以 $-1/2$，然后消去其他行：
  \[
    \left[
      \begin{array}{cccc|cccc}
        1 & 0 & 3 & 8 & -1 & 1 & 0 & 0 \\
        0 & 1 & -9 & -17 & 3 & -2 & 0 & 0 \\
        0 & 0 & 25 & 40 & -7 & 5 & 1 & 0 \\
        0 & 0 & -4 & -3 & 1 & -1 & 0 & 1
    \end{array}\right]
  \]

  \textbf{第 3 步：处理第 3 列}

  主元为 $25$，将第 3 行除以 $25$，然后消去其他行：
  \[
    \left[
      \begin{array}{cccc|cccc}
        1 & 0 & 0 & 16/5 & -4/25 & 2/5 & -3/25 & 0 \\
        0 & 1 & 0 & -13/5 & 12/25 & -1/5 & 9/25 & 0 \\
        0 & 0 & 1 & 8/5 & -7/25 & 1/5 & 1/25 & 0 \\
        0 & 0 & 0 & 17/5 & -3/25 & -1/5 & 4/25 & 1
    \end{array}\right]
  \]

  \textbf{第 4 步：处理第 4 列}

  主元为 $17/5$，将第 4 行除以 $17/5$，然后消去其他行，得到化简行阶梯形 $[I | A^{-1}]$：
  \[
    \left[
      \begin{array}{cccc|cccc}
        1 & 0 & 0 & 0 & -4/85 & 10/17 & -23/85 & -16/17 \\
        0 & 1 & 0 & 0 & 33/85 & -6/17 & 41/85 & 13/17 \\
        0 & 0 & 1 & 0 & -19/85 & 5/17 & -3/85 & -8/17 \\
        0 & 0 & 0 & 1 & -3/85 & -1/17 & 4/85 & 5/17
    \end{array}\right]
  \]

  \textbf{逆矩阵}

  \[
    A^{-1} =
    \begin{pmatrix}
      -\dfrac{4}{85} & \dfrac{10}{17} & -\dfrac{23}{85} &
      -\dfrac{16}{17} \\[6pt]
      \dfrac{33}{85} & -\dfrac{6}{17} & \dfrac{41}{85} & \dfrac{13}{17} \\[6pt]
      -\dfrac{19}{85} & \dfrac{5}{17} & -\dfrac{3}{85} & -\dfrac{8}{17} \\[6pt]
      -\dfrac{3}{85} & -\dfrac{1}{17} & \dfrac{4}{85} & \dfrac{5}{17}
    \end{pmatrix}
  \]

  验证可得 $A \cdot A^{-1} = I$，即逆矩阵正确。
\end{solution}

\section{上机习题 2}

先用你所熟悉的计算机语言将平方根法和改进的平方根法编写成求解程序，然后用你编写的程序分别求解下列方程组。

\subsection*{(1) 给定矩阵}

100 阶三对角矩阵：

\[
  \begin{bmatrix}
    10 & 1 & & & & &\\
    1 & 10 & 1 & & & & \\
    & 1 & 10 & 1 & & & \\
    & & \ddots & \ddots & \ddots & & \\
    & & & 1 & 10 & 1 & \\
    & & & & 1 & 10 & 1
  \end{bmatrix}
\]

$b$ 随机

\subsection*{(2) Hilbert 矩阵}

系数矩阵 $A$ 为 40 阶 Hilbert 矩阵，即系数矩阵 $A$ 的第 $i$ 行第 $j$ 列元素为
\[
  a_{ij} = \frac{1}{i+j-1}
\]
向量 $\vec{b}$ 的第 $i$ 个分量为
\[
  b_i = \sum_{j=1}^{n} \frac{1}{i+j-1}
\]

\begin{solution}
  \mbox{}
  \lstinputlisting[language=python,caption={Cholesky 实现}]{cholesky.py}
\end{solution}

\section{上机习题 3}

用第 1 题（第一讲的问题）的程序求解第 2 题的两个方程组，并比较两种求解方程的优缺点。

\begin{solution}
  \mbox{}
  \lstinputlisting[language=python,caption={测试代码}]{exercise2.py}
  \lstinputlisting[language={},caption={输出结果}]{output.txt}

  这里可以看出，朴素的 Cholesky 分解因为处理平方根的问题，误差较大，在 Hilbert
  矩阵的情况下甚至出现数值不稳定的现象，导致无法求解。而改进的 Cholesky 分解则没有这个问题，能够正确求解出两个方程组的解。

  另一方面，在性能方面，改进的 Cholesky 分解在处理大规模矩阵时性能表现更好，在随机数据的压力测试下相比列主元的高斯消元法有明显的优势。
\end{solution}

\end{document}
